\clearpage
\subsection{RQF-03 — Anti-spam y validación de email}

\renewcommand{\arraystretch}{1.3}

% --- 1. Información del Cliente ---
\noindent
\begin{tabular}{|p{5cm}|p{10.5cm}|}
\hline
\rowcolor{secVerde}\multicolumn{2}{|l|}{\color{white}\textbf{1. INFORMACIÓN DEL CLIENTE}} \\
\hline
Nombre de la empresa  & UTEQ \\
\hline
Nombre del Contacto   & Angélica Garduño Bustamante \\
\hline
Fecha de Elaboración  & 01/06/2026 \\
\hline
\end{tabular}

\vspace{0.3cm}

% --- 2. Información General ---
\noindent
\begin{tabular}{|p{5cm}|p{10.5cm}|}
\hline
\rowcolor{secVerde}\multicolumn{2}{|l|}{\color{white}\textbf{2. INFORMACIÓN GENERAL}} \\
\hline
Nombre del Requerimiento     & RQF 03. Anti-spam y validación de email \\
\hline
Tipo de Requerimiento        & Funcional \\
\hline
Fecha requerida de solución  & 30/06/2026 \\
\hline
Prioridad                    & Alta \\
\hline
\end{tabular}

\vspace{0.3cm}

% --- 3. Descripción del Requerimiento ---
\noindent
\begin{tabular}{|p{4.5cm}|p{11cm}|}
\hline
\rowcolor{secVerde}\multicolumn{2}{|l|}{\color{white}\textbf{3. DESCRIPCIÓN DEL REQUERIMIENTO}} \\
\hline
Objetivo del Requerimiento &
Implementar un mecanismo de validación de email mediante expresión regular y un sistema anti-spam que deshabilite los botones de los formularios durante las peticiones HTTP, evitando envíos duplicados al backend. \\
\hline
Beneficio esperado &
Reducción de la carga al servidor por peticiones duplicadas, mejora de la experiencia de captura de datos del usuario y prevención de errores de validación previos al envío. \\
\hline
Alcance &
- Aplicar validación de email en \textit{login}, \textit{register} y \textit{forgot-password}. \par
- Deshabilitar el botón de envío durante la petición HTTP. \par
- Mostrar borde rojo y mensaje inline en inputs inválidos. \par
- No modifica la lógica del backend. \\
\hline
Ubicación &
\url{https://github.com/MarkG777/AbejaNet} \\
\hline
Descripción de Usuarios y Responsables &
\textbf{Usuarios:} \par
\textbf{Usuario (Cliente):} interactúa con los formularios de autenticación de la app móvil. \par
\vspace{0.1cm}
\textbf{Responsables:} \par
\textbf{App Móvil (React Native):} ejecuta la validación local antes de enviar la petición. \par
\textbf{API Backend:} no requiere cambios; solo recibe peticiones bien formadas. \\
\hline
Descripción funcional &
1. El usuario ingresa al formulario de inicio de sesión, registro o recuperación de contraseña. \par
2. Al escribir el email, el sistema valida el formato en tiempo real con una expresión regular. \par
3. Si el formato es inválido, el input muestra borde rojo y mensaje inline; el botón permanece deshabilitado. \par
4. Cuando los datos son válidos y el usuario presiona el botón de envío, el botón se deshabilita inmediatamente y muestra un indicador de carga. \par
5. Mientras la petición está en curso, el formulario no acepta nuevas interacciones. \par
6. Al recibirse la respuesta del backend, el botón se reactiva y se muestra el resultado correspondiente. \\
\hline
Restricciones &
- La validación local no sustituye la validación del backend. \par
- Debe mantenerse la compatibilidad con el sistema de autenticación JWT existente. \\
\hline
Precedencias & RQF 02. Registro de clientes \\
\hline
Documentación requerida & Manuales de usuario actualizados con la nueva retroalimentación visual. \\
\hline
Observaciones &
Se reutiliza la librería de iconos \texttt{Ionicons} ya presente en el proyecto para los indicadores visuales. \\
\hline
\end{tabular}

\vspace{0.3cm}

% --- 4. Condiciones de Prueba ---
\noindent
\begin{tabular}{|p{5.1cm}|p{5.1cm}|p{5.1cm}|}
\hline
\rowcolor{secVerde}\multicolumn{3}{|l|}{\color{white}\textbf{4. CONDICIONES DE PRUEBA}} \\
\hline
\rowcolor{cellGris}\textbf{Condición} & \textbf{Validación} & \textbf{Resultado} \\
\hline
El usuario ingresa un email con formato inválido en la pantalla de inicio de sesión. &
El sistema debe detectar el formato incorrecto antes de enviar la petición y mostrar un mensaje inline. &
El input se marca con borde rojo y el botón de envío permanece deshabilitado. \\
\hline
El usuario presiona repetidamente el botón \textit{``Iniciar sesión''} mientras la petición está en curso. &
El sistema debe deshabilitar el botón al primer clic para evitar envíos duplicados. &
Únicamente se realiza una petición HTTP al backend, sin importar el número de clics. \\
\hline
\end{tabular}
